\documentclass[border=0pt]{standalone}
% The shared preamble of every docs diagram (docs/diagrams/<name>.tex).
% dvisvgm's driver first, so patterns and opacity survive into the SVG.
\def\pgfsysdriver{pgfsys-dvisvgm.def}
\usepackage[siunitx,RPvoltages]{circuitikz}
\usepackage{xcolor}
% Latin Modern: outline fonts for every glyph a label uses, µ included.
\usepackage[T1]{fontenc}
\usepackage{lmodern}
\usetikzlibrary{arrows.meta,backgrounds,calc,patterns,positioning}
\tikzset{
  % A white card, so diagrams read in mdBook's dark themes too.
  every picture/.append style={show background rectangle,
    background rectangle/.style={fill=white,inner sep=4pt}},
  % Layout geometry: layers told apart by fill and hatching, never colour alone.
  layer a/.style={draw=black,line width=0.6pt,fill=gray!20},
  layer b/.style={draw=black,line width=0.6pt,fill=gray!45},
  layer c/.style={draw=black,line width=0.6pt,fill=white,postaction={pattern=north east lines}},
  layer d/.style={draw=black,line width=0.6pt,fill=white,postaction={pattern=crosshatch dots}},
  marker/.style={draw=black,dashed,line width=0.6pt},
  % A measured distance and its label.
  dim/.style={{Latex[length=2mm]}-{Latex[length=2mm]},line width=0.5pt},
  dim label/.style={font=\small,fill=white,inner sep=1pt},
  % Where a rule fails: thick, dashed and labelled, not only red.
  violation/.style={draw=red!75!black,line width=1.4pt,dash pattern=on 3pt off 2pt},
  violation label/.style={font=\small\bfseries,text=red!75!black},
  caption/.style={font=\small},
}

\begin{document}
\begin{circuitikz}
  % The inverter at x = 0 between VDD (y 2.6) and VSS (y -2.6).
  \draw (0,1) node[pmos,anchor=D](ap){} (0,-1) node[nmos,anchor=D](an){};
  \draw (ap.D) -- (an.D);
  \draw (ap.G) -- (an.G) coordinate[pos=0.5](ain);
  \draw (ap.S) -- (ap.S |- 0,2.6);
  \draw (an.S) -- (an.S |- 0,-2.6);
  \draw (-1.6,2.6) -- (1.8,2.6) node[right]{VDD, 1.8 V};
  \draw (-1.6,-2.6) -- (1.8,-2.6) node[right]{VSS, 0 V};
  \draw (0,0) to[short,*-] (1.2,0) node[right]{Y: 0 V};
  % The switch from VDDH to X, gate EN.
  \draw (-3.4,3.0) node[nmos,anchor=D](s){};
  \draw (s.D) -- ++(0,0.3) node[above]{VDDH, 3.3 V};
  \draw (s.G) -- ++(-0.4,0) node[left,align=right]{EN: input,\\may be anything};
  \draw (s.S) |- (ain);
  \node[caption,anchor=south] at (-2.3,0.05) {X: 3.3 V};
  \node[caption,anchor=west] at (-3.3,2.2) {conducts};
  % The failing devices.
  \draw[violation] (-1.15,-2.1) rectangle (0.5,2.1);
  \node[violation label,anchor=north] at (-0.3,-2.75) {3.3 V on a 1.95 V oxide};
  \node[caption,anchor=north,align=center] at (-0.3,-3.25) {EN tied to VSS instead: the switch is off, X is reached\\by no supply, and the inverter is not examined};
\end{circuitikz}
\end{document}
